\section{Introduction}
\label{sec:intro}

Diabetes now affects roughly one adult in ten. The International Diabetes Federation estimated that 536.6 million people aged 20--79, or 10.5\% of that age group, were living with the disease in 2021, and it expects the figure to reach 783.2 million (12.2\%) by 2045~\cite{idf2022}. Worse, many people with type~2 diabetes and most with prediabetes do not know they have it, so complications build unnoticed and much effort goes into deciding who should be tested first.

This question matters to me because of the work I want to do. My mission is to revolutionise healthcare in sub-Saharan Africa through technology, creating accessible, equitable and scalable solutions for underserved rural communities, and it is in exactly those communities that the diagnosis gap is widest. The IDF reports that 24 million African adults had diabetes in 2021 and that 54\% of them were undiagnosed~\cite{idfafrica2021}. Kenya's 2015 national STEPS survey found that more than half of adults with diabetes were unaware of their condition, and measured prevalence was lower in rural areas~\cite{mohamed2018}, where fewer people are ever tested. In a rural clinic without a laboratory, a short questionnaire administered by a community health worker is about the most scalable screening tool available. Before anyone builds such a tool, though, someone has to establish how much a questionnaire can actually reveal, and what it costs in missed cases and false alarms. That is the gap this project tries to fill.

I could not do this with African data. No openly available survey from the region yet combines a large sample with richly documented predictors; the Kenyan STEPS analysis, for example, covered 4,069 adults~\cite{mohamed2018}. Instead I use the US Behavioral Risk Factor Surveillance System (BRFSS), which interviews hundreds of thousands of adults each year about chronic conditions, behaviour, access to care and demographics~\cite{cdc2015data}, and whose 2015 file is documented variable by variable in its codebook~\cite{cdc2015codebook}. I treat it as a testbed. Whatever it teaches has to be revalidated on local data before it is used anywhere in Africa. The BRFSS has also attracted a growing machine learning literature, but as Section~\ref{sec:lit} shows, much of that work simplifies the outcome to a binary label, uses a cleaned subset of the data, and reports accuracy at default thresholds, choices that make the task look easier than it really is.

My research question is therefore: \emph{How effectively can traditional machine learning and deep learning models classify diabetes status from population health indicators, and how do their predictive performance, interpretability, computational efficiency and error patterns differ?} To answer it, I kept the survey's own three categories of no diabetes, prediabetes and diabetes, built a leakage-safe pipeline, and compared logistic regression, random forests and gradient boosting with Sequential and Functional TensorFlow/Keras networks across 17 tracked experiments, before examining the selected models' thresholds, errors, explanations and subgroup performance.

Going in, I expected gradient boosting to give the best discrimination, neural networks to match it without beating it, and logistic regression to trail only slightly. Benchmark studies point that way for tabular data~\cite{grinsztajn2022,shwartzziv2022}, clinical reviews find little advantage for complex models~\cite{christodoulou2019}, and most survey answers are simple yes/no or ordered categories whose effects are likely to be additive. I also expected prediabetes to be very hard to separate from its neighbours. Throughout, I keep in mind that the labels are self-reported answers to the question of whether a doctor ever said the person had diabetes. They are not laboratory results, and none of the models here is a diagnostic tool.
